\documentclass[aspectratio=169,10pt]{beamer}
\usepackage[T1]{fontenc}
\usepackage{lmodern}
\usepackage{listings}
\usepackage{xcolor}

\definecolor{Ink}{HTML}{172B3A}
\definecolor{Blue}{HTML}{235789}
\definecolor{Muted}{HTML}{526574}
\definecolor{Paper}{HTML}{F5F7F8}
\definecolor{Code}{HTML}{EAF0F4}
\setbeamercolor{normal text}{fg=Ink,bg=white}
\setbeamercolor{frametitle}{fg=Ink,bg=white}
\setbeamercolor{title}{fg=Ink}
\setbeamercolor{structure}{fg=Blue}
\setbeamercolor{block title}{fg=Blue,bg=white}
\setbeamercolor{block body}{fg=Ink,bg=Paper}
\setbeamertemplate{navigation symbols}{}
\setbeamertemplate{footline}[frame number]
\setbeamersize{text margin left=0.7cm,text margin right=0.7cm}
\setlength{\parskip}{0.35em}

\lstset{
  basicstyle=\ttfamily\small,
  columns=fullflexible,
  keepspaces=true,
  showstringspaces=false,
  breaklines=true,
  backgroundcolor=\color{Code},
  frame=single,
  rulecolor=\color{Code},
  xleftmargin=0.4em,
  xrightmargin=0.4em,
  aboveskip=0.6em,
  belowskip=0.6em
}

\title{Turning LAMBDA into Python}
\subtitle{A register-based compiler, using recursive factorial}
\author{CS413}
\date{Fall 2026}

\begin{document}

\begin{frame}[plain]
  \titlepage
\end{frame}

\begin{frame}{Compiler pipeline}
  \begin{enumerate}
    \item \textbf{Parse source.} \texttt{d0exp\_parse} turns LAMBDA text into expression nodes.
    \item \textbf{Check types.} The type checker can reject invalid expressions before compilation.
    \item \textbf{Compile the AST.} \texttt{d0exp\_compile} lowers each node to registers and Python statements.
    \item \textbf{Run Python.} The generated module defines \texttt{\_lambda\_main} and stores its result in \texttt{\_lambda\_result}.
  \end{enumerate}
  \vspace{0.6em}
  \begin{block}{Important interface detail}
    The compiler accepts an AST, not source text. A typical caller parses first, then passes the AST to \texttt{d0exp\_compile}.
  \end{block}
\end{frame}

\begin{frame}[fragile]{Input program: factorial}
  \begin{lstlisting}[language=Python]
let
  fact : int -> int =
    fix fact(n: int): int =>
      if n <= 1 then 1 else n * fact(n - 1)
in
  (fact(0), fact(1), fact(5), fact(6))
end
  \end{lstlisting}
  \vspace{0.3em}
  \begin{itemize}
    \item \texttt{fix} binds both the recursive function name and its parameter.
    \item The condition selects one branch; the other branch must not run.
    \item The final tuple makes several results visible in one evaluation.
  \end{itemize}
\end{frame}

\begin{frame}[fragile]{The parser produces an AST}
  The outer expression is a sequential binding followed by a tuple result:
  \begin{lstlisting}
D0Elet(some name elided)
  value: D0Efix(fact, n, body)
    body: D0Eif0
      test: D0Eop2("<=", D0Evar(n), D0Eint(1))
      then: D0Eint(1)
      else: D0Eop2("*", D0Evar(n), recursive_call)
  body: D0Etupl(fact(0), fact(1), fact(5), fact(6))
  \end{lstlisting}
  \vspace{0.3em}
  \texttt{D0Eint}, \texttt{D0Efix}, and \texttt{D0Eif0} are dataclass variants. Types are annotations on AST nodes and are erased by code generation.
\end{frame}

\begin{frame}[fragile]{AST classes describe LAMBDA syntax}
  Expression variants inherit from \texttt{D0E000}; declarations inherit from
  \texttt{D0C000}. Their dataclass fields hold child nodes and source names.
  \begin{columns}[T]
    \column{.48\textwidth}
    \begin{lstlisting}
D0Eint(value)
D0Evar(name)
D0Eop1(operator, operand)
D0Eop2(operator, left, right)
D0Eapp(function, argument)
D0Eif0(test, then, otherwise)
    \end{lstlisting}
    \column{.48\textwidth}
    \begin{lstlisting}
D0Elam(parameter, body, annotation)
D0Efix(name, parameter, body, ...)
D0Elet(name, value, body)
D0Elets(declarations, body)
D0Etupl(items)
D0Eproj(tuple, zero_based_index)
    \end{lstlisting}
  \end{columns}
  \smallskip
  \texttt{D0Cval} represents a named declaration, including its optional type
  annotation and explicit type parameters. The aliases \texttt{d0exp} and
  \texttt{d0cls} name the expression and declaration base types.
\end{frame}

\begin{frame}[fragile]{Type classes represent checker data}
  \begin{columns}[T]
    \column{.48\textwidth}
    \begin{lstlisting}
TInt()        TBool()
TTuple(items)
TFun(argument, result)
TList(element) TOption(element)
TMeta(ident)  TPoly(name)
    \end{lstlisting}
    The variants form the \texttt{simple\_type} union. The type dataclasses
    are frozen, so a type node is immutable after construction.
    \column{.48\textwidth}
    \begin{lstlisting}
_Scheme(
  quantified={...},
  typ=TFun(TMeta(0), TMeta(0)))
    \end{lstlisting}
    \texttt{TMeta} holds an inference variable. \texttt{TPoly} marks a declared
    type parameter. \texttt{\_Scheme} pairs a type with the inference
    variables quantified by a polymorphic declaration.
  \end{columns}
\end{frame}

\begin{frame}[fragile]{Runtime values and environments}
  The interpreter has its own class hierarchy; evaluating an AST produces
  \texttt{D0V} values rather than Python primitives.
  \begin{columns}[T]
    \column{.48\textwidth}
    \begin{lstlisting}
D0Vint(number)   D0Vbtf(flag)
D0Vlist(items)  D0Voptn(option)
D0Vtupl(items)
D0Vlam(environment, lambda_ast)
D0Vfix(environment, fix_ast)
    \end{lstlisting}
    A function value keeps the syntax it will run and the environment needed
    for its free variables. This pair is a closure.
    \column{.48\textwidth}
    \begin{lstlisting}
ENVnil()
ENVcns(name, value, rest)
    \end{lstlisting}
    Environments are linked bindings. Lookup checks the newest binding first,
    so an inner binding naturally shadows an outer one. The type \texttt{d0val}
    aliases the runtime value base class.
  \end{columns}
\end{frame}

\begin{frame}[fragile]{Classes supporting parsing and compilation}
  \begin{itemize}
    \item \texttt{\_Token} stores a token's kind, text, and source offset;
      \texttt{\_Parser} consumes tokens and builds the AST.
    \item \texttt{LambdaSyntaxError} reports malformed source with line and
      column; \texttt{TypeCheckError} reports static type failures.
    \item \texttt{\_TypeSolver} creates \texttt{TMeta} variables and records
      substitutions while unifying types.
    \item \texttt{\_PyCompiler} walks the expression AST;
      \texttt{\_PyComp} carries the generated result name and instruction list.
    \item Instructions are dataclass variants: \texttt{\_PyAssign},
      \texttt{\_PyDef}, and \texttt{\_PyIf}.
  \end{itemize}
  These helper classes support separate stages. The expression and type nodes
  are the shared data passed between parsing, checking, interpreting, and
  Python emission.
\end{frame}

\begin{frame}[fragile]{Compilation uses result registers}
  Each recursive compiler call returns a result name and ordered instruction
  objects. The emitter turns these objects into Python statements:
  \begin{lstlisting}
_PyComp(result="_v6", instructions=[
    _PyAssign(target="_v4", value="1"),
    _PyAssign(target="_v5", value="(_a1 - _v4)"),
    _PyAssign(target="_v6", value="_v1(_v5)"),
])
  \end{lstlisting}
  \begin{itemize}
    \item \texttt{reg\_new()} creates names such as \texttt{\_v1}, \texttt{\_v2}.
    \item The compile environment maps LAMBDA names to Python names, for example \texttt{n -> \_a1}.
    \item \texttt{\_PyDef} and \texttt{\_PyIf} contain nested instruction lists, so function and branch bodies stay structured.
  \end{itemize}
\end{frame}

\begin{frame}[fragile]{Instruction objects and Python emission}
  Assignment instructions store a target and expression. Control-flow
  instructions store child blocks and their result registers:
  \begin{lstlisting}
_PyIf(test="_v3", result="_v9",
  then_body=[_PyAssign("_v4", "1")], then_result="_v4",
  else_body=[_PyAssign("_v5", "2")], else_result="_v5")
  \end{lstlisting}
  \begin{itemize}
    \item The recursive emitter prints each instruction at the requested indentation.
    \item \texttt{d0exp\_compile(ast)} returns the complete generated program.
    \item \texttt{d0exp\_emit(ast)} prints that program to standard output.
  \end{itemize}
\end{frame}

\begin{frame}[fragile]{Scope and evaluation order}
  The compiler environment maps each LAMBDA name to a generated Python local.
  Extending a scope copies that map, so shadowing affects only the body.
  \begin{columns}[T]
    \column{.48\textwidth}
    \textbf{A let aliases its value register}
    \begin{lstlisting}
value = compile(rhs, env)
body_env = dict(env)
body_env["x"] = value.result
body = compile(body, body_env)
    \end{lstlisting}
    No extra Python assignment is needed. Sequential declarations extend the
    environment after each initializer.
    \column{.48\textwidth}
    \textbf{Instructions preserve call-by-value order}
    \begin{lstlisting}
function instructions
argument instructions
result = function(argument)
    \end{lstlisting}
    Binary operators likewise emit left operand code, then right operand code,
    followed by the operation.
  \end{columns}
\end{frame}

\begin{frame}[fragile]{Recursive functions become Python functions}
  A LAMBDA \texttt{fix fact(n) => body} becomes a Python \texttt{def}. The compiler maps the recursive name to the function's own generated name.
  \begin{lstlisting}
def _v1(_a1):
    _v2 = 1
    _v3 = (_a1 <= _v2)
    if _v3:
        _v4 = 1
        _v9 = _v4
    else:
        _v5 = 1
        _v6 = (_a1 - _v5)
        _v7 = _v1(_v6)       # recursive call
        _v8 = (_a1 * _v7)
        _v9 = _v8
    return _v9
  \end{lstlisting}
  The nested \texttt{def} also closes over its surrounding registers.
\end{frame}

\begin{frame}[fragile]{Conditionals preserve LAMBDA evaluation}
  The compiler emits branch bodies inside Python's \texttt{if}/\texttt{else}. Both branches are compiled, but only the selected branch executes.
  \begin{lstlisting}
if _v3:
    _v4 = 1
    _v9 = _v4
else:
    _v5 = 1
    _v6 = (_a1 - _v5)
    _v7 = _v1(_v6)
    _v8 = (_a1 * _v7)
    _v9 = _v8
  \end{lstlisting}
  \texttt{\_v9} is the merge register: each branch assigns it, and the surrounding code uses it as the conditional's result.
\end{frame}

\begin{frame}[fragile]{Generated module and running the example}
  The compiler wraps the instruction list in a zero-argument entry point, adds runtime helpers for LAMBDA lists, options, and printing, then invokes the entry point.
  \begin{lstlisting}
def _lambda_main():
    # generated instructions
    return _v18

_lambda_result = _lambda_main()
  \end{lstlisting}
  From \texttt{LAMBDA2/mytest}:
  \begin{lstlisting}[language=Python]
make check           # type-check and compare with the interpreter
make show-factorial  # print the generated Python source
  \end{lstlisting}
  The factorial example returns \texttt{(1, 1, 120, 720)}.
\end{frame}

\begin{frame}[fragile]{Runtime values and helper calls}
  Generated code uses Python's built-in values where they fit the LAMBDA data
  model, and helper functions for operations that need checks or special
  representations.
  \begin{columns}[T]
    \column{.48\textwidth}
    \begin{itemize}
      \item Integers and booleans use Python \texttt{int} and \texttt{bool}.
      \item Lists and tuples use Python lists and tuples.
      \item \texttt{list\_head}, \texttt{list\_tail}, and \texttt{list\_cons}
        call runtime helpers that check their inputs.
    \end{itemize}
    \column{.48\textwidth}
    \begin{itemize}
      \item \texttt{optn\_nil} is \texttt{None}; \texttt{optn\_cons(x)} is a
        tagged pair, so an option containing \texttt{None} stays distinct.
      \item Source \texttt{/} emits Python \texttt{//} for integer division.
      \item \texttt{print} formats the value, prints it, and returns unit
        (\texttt{()}).
    \end{itemize}
  \end{columns}
\end{frame}

\end{document}
